\section{Entropy-Aware Hashing}

This project was completed for CSC 519: Computer Networks.

\subsection{Background}

% relevance of hash functions
% limitations of current hash functions (general)
Hashing is ubiquitous in computing. For example, routing table lookups, intrusion detection systems, load distribution, software defined networking, network 
address translation and application-layer caches rely on hashing~\cite{kurose2025}. Hashing is already highly optimized, so improving per-byte throughput is 
difficult. However, by specializing hash functions to particular applications, increasing throughput might be possible. For example, hashing fewer bytes. \\


% entropy learned hashing, sede
% main idea(s): hash fewer bytes, use efficient instructions
% research questions:
% - ELH + sede do not use some information: dist of keys, data structure
% incorporate this
\noindent
In fact, \cite{mitzenmacher2008} showed that hash functions work by extracting entropy from the input data. Several works, including entropy learned 
hashing~\cite{hentschel2022} and \texttt{Sede}~\cite{hoffmann2025} exploit this by learning the position of high entropy bytes and only hashing them. However, 
their specialization is incomplete: they do not incorporate non-uniform key frequencies or the target data structure into the cost model. Therefore, we ask:
\begin{enumerate}
    \item How can hash functions be automatically specialized to a target distribution, data structure and architecture?
\end{enumerate}

\subsection{Summary}

\subsubsection{Entropy-Aware Hashing}

% overview
% - illustrate framework by optimizing
% - std unordered set on x64, URL key dataset
We define two degrees of freedom: hash function form and cost model. To specialize a hash function, we sweep hash function forms. For each form, we greedily 
select data chunks to hash, minimizing the cost of the specialized function. We select between specialized forms by benchmarking. We illustrate the framework by 
specializing a hash function for \texttt{std::unordered\_set} containing a dataset of URLs on the \texttt{x64} architecture. \\

% hash function forms:
% - hash function
% - alignment
% - data size
% datatype: byte, word, or double-word
% which hash functions: MurmurHash3, Thomas Wang's, Stafford's, XXHash, AES round, GXHash
% alignment: none or datatype aligned
\noindent
Unlike entropy learned hashing which selected single bytes, we allow our hash function forms to select either single bytes, words, or double-words. Selecting at 
higher granularity means lower entropy bytes might be selected, but may permit more efficient implementations. We also provide an option to force alignment to 
this data type size. Some architecture may have penalties for accessing unaligned data. Finally, the hash function form includes the mixing function selected 
from MurmurHash3, Thomas Wang's, Stafford's, XXHash, GXHash and AES round. \\

% cost models:
% - unweighted collision cost
% - weighted collision cost
% - maximum weighted probe cost
% - minimum weighted probe cost
\noindent
Entropy learned hashing selected a subset of bytes to minimize the unweighted collision cost (the expected number of collisions in the resulting hash table). 
We consider three additional cost models. The weighted collision cost considers the distribution of keys: collisions between high frequency keys are more 
expensive than low frequency keys. The maximum weighted probe cost is the expected number of probes in the worst possible resulting hash table. Likewise, the 
minimum weight probe cost is the expected number of probes in the best possible resulting hash table. Best and worst are defined with respect to the chain 
orders. \\

\noindent
We find that selecting unaligned double-words with unweighted collision cost and applying a single round of AES yields the fastest hash function and hashset
queries. Hash latency is reduced from 12.1ns for \texttt{std::hash} to 4.5ns. Query latency is reduced from 32.9ns for \texttt{std::hash} to 14.7ns. Build time
is not impacted. In general, we find that specializing to the data distribution and data structure does not provide a strong benefit, but specialization to the
architecture with unaligned double-words and AES instructions yields a significant speedup. 

% findings: 2x speedup, DS/dist specialization do not provide strong benefit
% double-word, unweighted collision cost, unaligned, single round AES
% hash latency: 12.1 (std::hash) -> 4.5 ns
% query latency: 32.9 (std::hash) -> 14.7 ns
% no change in build time

\subsubsection{Future Work}

We see five promising directions for future work. First, additional cost models, hash function forms, datasets, data structures, and architectures should be 
evaluated. Second, entropy-aware hashing should adapt to evolving data distributions automatically. This will require learning the position of high entropy 
regions online. Third, entropy-aware hashing should handle varying-length keys. This could be framed as a multiple sequence alignment problem. Fourth, the cost
model could be automatically learned from the dataset, data structure and architecture with deep learning. Fifth, the hardness of entropy maximization and the 
greedy heuristic should be formally analyzed. There might be better optimizers than greedy.

% prove hardness, greedy is approximation
% dynamicly learn entropy
% different length keys (sequence alignment)
% automatically learn cost model from DS, data and arch.

% TODO: A 400–500-word summary of the technical contributions of the work, what was
% accomplished overall, and the limitations of their work.

\subsection{Reflection}

% SIMD/instrisics
This project strengthened my ability to optimize software with SIMD instructions. This was my first experience implementing efficient algorithms with SIMD 
instructions. I learned to read the \texttt{x64} instruction set manual and optimization manual. I identified that a single round of AES would be a highly
efficient mixing function. \\

% approximation, hardness
\noindent
This project also developed my knowledge of theoretical computer science. I attempted to prove that selecting finding the subset with maximum entropy was 
\textbf{NP-Hard}, and that the greedy heuristic was an optimal approximation algorithm. Although I did not succeed, I learned a lot about approximation 
algorithms.  


% TODO: A 250-word self reflection of what skills were developed or enhanced in completing those
% projects.